\documentclass[border=1pt]{standalone}

\usepackage[T1]{fontenc}
\usepackage[utf8]{inputenc}
\usepackage{lmodern}
\usepackage{microtype}
\usepackage{graphicx}
\usepackage{booktabs}
\usepackage{amsmath,amssymb,amsfonts}
\usepackage{bm}
\usepackage{siunitx}



\usepackage{tikz}
\usepackage{pgfplots}
\usepackage{pgfplotstable}
\usepackage{caption}
\pgfplotsset{compat=newest}




\newcommand{\papertitle}{Spectroscopic Constraints on Kelvin-Like Internal Modes in Effective Filament Models of Atomic Bound States}
\newcommand{\papershorttitle}{Kelvin-Like Internal Modes in Atomic Bound States}
\newcommand{\paperdoi}{10.5281/zenodo.18018084}
\newcommand{\paperorcid}{0009-0006-1686-3961}
\newcommand{\paperemail}{info@omariskandarani.com}
\newcommand{\papercity}{Groningen}
\newcommand{\papercountry}{The Netherlands}
\newcommand{\paperkeywords}{Kelvin modes, effective filament models, atomic spectroscopy, excitation gap, low-energy reduction}

\raggedbottom
\\pgfplotsset{compat=1.18}


\providecommand{\papertitle}{Spectroscopic Constraints on Kelvin-Like Internal Modes in Effective Filament Models of Atomic Bound States}
\providecommand{\papershorttitle}{Kelvin-Like Internal Modes in Atomic Bound States}
\providecommand{\paperdoi}{10.5281/zenodo.18018084}
\providecommand{\paperorcid}{0009-0006-1686-3961}
\providecommand{\paperemail}{info@omariskandarani.com}
\providecommand{\papercity}{Groningen}
\providecommand{\papercountry}{The Netherlands}
\providecommand{\paperkeywords}{Kelvin modes, effective filament models, atomic spectroscopy, excitation gap, low-energy reduction}
\providecommand{\epsilonK}{\varepsilon_K}

\usetikzlibrary{
    spath3,
    intersections,
    arrows,
    arrows.meta,
    knots,
    calc,
    hobby,
    decorations.pathreplacing,
    shapes.geometric,
}

% ---------------- Global styles (safe for 'knots') ----------------
\tikzset{
    knot diagram/every strand/.append style={
        line cap=round,
        line join=round,
        ultra thick,
        black
    },
    every knot/.style={line cap=round,line join=round,very thick},
    strand/.style={line cap=round,line join=round,line width=3pt,draw=black},
    over/.style={preaction={draw=white,line width=6.5pt}},
% DO NOT set a global "every path" here; it breaks internal clip paths.
}
\providecommand{\swirlarrow}{\rightsquigarrow}
\begin{document}
\begin{tikzpicture}
[
                width=6cm,
                height=4.5cm,
                title={Kelvin Excitation Spectrum},
                title style={font=\bfseries\tiny},
                xlabel={$k_m$ (mode index)},
                ylabel={$E_m$ [eV]},
                axis lines=left,
                ticks=none,
                xmin=0, xmax=6,
                ymin=0, ymax=800,
                clip=false
                ]
                \pgfmathsetmacro{\DeltaK}{500} 
                \fill[gray!20] (axis cs:0,0) rectangle (axis cs:6,\DeltaK);
                \draw[dashed, gray!60] (axis cs:0,\DeltaK) -- (axis cs:6,\DeltaK);
                \node[align=center,font=\tiny] at (axis cs:3,250) {Forbidden\\ region};
                \foreach \x in {1,2,3,4,5} {
                    \addplot[only marks, mark=*, red, mark size=2pt] coordinates {(\x, {\DeltaK + 40*\x})};
                }
                \node[anchor=west,font=\tiny] at (axis cs:6,\DeltaK+10) {$\Delta_K$};
\end{tikzpicture}
\end{document}